\RequirePackage[T1]{fontenc}
\documentclass[conference,10pt]{IEEEtran}
\usepackage{amsmath,amssymb,graphicx,url}
\title{Current-STA Transactional Repair for Timing--Resource Tradeoffs in 3D FPGA Routing}
\author{\IEEEauthorblockN{Anonymous development manuscript}}
\begin{document}
\maketitle
\begin{abstract}
Legal FPGA routes can occupy different points in a timing--resource design space. In a fixed 3D routing resource graph, sharing a vertical resource may save occupancy but restrict a critical connection or create capacity competition. We study a current-STA repair layer that selects complete connection groups, performs native searches on private route-tree copies, and commits only after physical, capacity and native timing checks. Retained-crossing domains encourage sharing; a failed guided search can release only its network's domain while preserving successful restorations elsewhere. Development results include three automatic timing repairs of a strong native BNN route, reducing critical-path delay from13.607436 to13.228146 ns with forty additional planar wire units. Subsequent resource refinement exposes distinct low-wire and low-vertical-occupancy alternatives at the same delay. On eltwise, a uniform K8 grouping cap reduces actual SB occupancy to24 and CB occupancy to965 at unchanged delay; an R1-based SPMV repair remains dominated by strong A0. Matched native controls distinguish selection, grouping and domain contributions. Prototype accounting exposes construction and validation overhead; complete cold algorithm timing and broader fixed-policy confirmation remain open.
\end{abstract}
\begin{IEEEkeywords}
3D FPGA, detailed routing, static timing analysis, transactional repair, multiobjective optimization
\end{IEEEkeywords}

\section{Problem and Motivation}
Three-dimensional FPGA fabrics couple directed reachability, scarce vertical channels and detailed timing. A different route can improve timing, save resources or exchange one benefit for another even when packing, placement and architecture remain fixed. A repair that does not reduce the final critical-path delay (CPD) may still be useful under an explicit resource or timing requirement. Conversely, a locally shorter connection is not evidence of a better complete route: another endpoint can become critical, or the replacement may prevent a different network from using a scarce channel.

VTR provides the detailed routing and timing infrastructure used here~\cite{vtr9paper}. The placement and architectural context follows the separation of design tradeoffs emphasized by Escaping Flatland~\cite{flatland}; our comparisons hold each physical input fixed. This work concerns post-legal routing decisions. It neither changes the physical timing model nor infers silicon area or power from reduced occupancy on a fixed device.

Initial shared-crossing guide experiments motivated the change of intervention point. Candidate coverage, search adoption, retained branches and native convergence can disconnect a guide's proxy score from final timing. We therefore operate on the current routed state and native STA, and evaluate an actual replacement rather than assuming that a different plan is adopted. The repair search itself remains native. The contributions under investigation are group selection, dependency scope, retained-domain decisions, acceptance and scheduling.

Our evidence distinguishes three questions: does an actual legal route improve a joint quality coordinate; which mechanism caused the change under matched local controls; and what does a complete fixed configuration cost? The first two have development evidence. The last requires further controlled full-flow measurements, rather than substituting a repair core or a shared research workflow.

\section{Objective and State Identity}
An input identity fixes the directed resource graph, packed netlist, placement, endpoint mapping, SDC and clock domains. A legal route $R$ uses actual graph edges and respects node capacities and native routing constraints. We record
\begin{equation}
 q(R)=(D(R),W(R),S(R),C(R),T(R)),
\end{equation}
where $D$ is native final CPD, $W$ is planar wire in the auditor's native resource units, $S$ is occupied SB pairs, $C$ is occupied CB output sources, and $T$ is comparable complete algorithm time. Paired SB half-nodes are one vertical resource; multiple CB edges from one output are not independent outputs. Guide recommendations are not occupancy measurements.

A final legal candidate is retained when it is not dominated on the measured, comparable coordinates. Quality projections use the first four coordinates and preserve uncertainty from printed CPD. They do not establish five-dimensional dominance when complete time is unknown. Different convergence rules, repair opportunities and implementation versions remain explicit, even for a descriptive comparison on the same physical inputs.

Final retention is separate from a branch's transition rule. We retain a strict CPD mode, $D(R')<D(R)$, and a fixed-resource mode that reduces CPD within registered absolute resource limits. A resource-under-timing mode permits wire tradeoffs within a fixed envelope while requiring nonincreasing selected resource objectives and at least one strict reduction. The current resource probes select $(S,C)$ as those objectives, cap CPD at the root value, cap wire at the root value plus0.1\%, and prohibit increases beyond root vertical counts. These are design parameters, not general termination conditions.

An endpoint-vector control fixes all endpoint/launch/capture identities and compares complete native setup slack values. It accepts only nonworse CPD, no setup-component regression and at least one strict setup improvement, with no increased hold violation. Native numeric values are compared without an epsilon. Resource policies impose their stated CPD and hold conditions; they do not silently inherit the stricter endpoint-vector requirement. Anchors do not roll forward, so a tolerance cannot accumulate timing degradation.

\section{Mechanism: Current-State Whole Transactions}
\subsection{Select targets and dependency scopes}
Timing selection follows the current native winning-path structure, including tied fan-ins, rather than a permanently retained list of endpoints. Capacity queries distinguish a promising extension from one blocked by other networks. Owner scopes and grouped targets can release dependencies that a single-connection repair leaves frozen.

Resource selection scans the actual current route trees. A candidate removes all sinks below a selected crossing subtree, preserves paths to the other sinks, and requires both structural release potential and a retained crossing alternative. The screen uses current slack margin; it is not a proof of final resource savings. Ranking is deterministic: released crossings per removed sink, fewer sinks, larger current margin, then network and pin IDs. A K8 group contains at most eight distinct networks; fewer are used when fewer qualify. Removed connections are restored in descending current raw criticality with deterministic ties. Accepted states receive a fresh current selector.

\subsection{Native searches and retained-crossing domains}
For each group we copy the selected trees and prune exactly the registered sink sets. Occupancy is updated as every local replacement is installed, so later searches see earlier resource use. Native weighted search uses the full device search box and the existing timing/congestion objective. Capacity-infeasible results are not committed. Real switch delays and final STA remain unchanged.

The domain control initially allows cross-layer edges only through crossing sources retained after pruning. It preserves useful sharing but can become infeasible under current capacities. The unrestricted native control receives the same current source, sink set, tree scope, objective and search range. Paired development probes evaluate both variants and use the registered domain-first acceptance order; native-only trajectories instead reselect from their own accepted unrestricted states. Those trajectories must not borrow a guided successor after divergence.

Automatic failure release makes the domain a revisable decision. If guided native search fails to return a capacity-feasible branch, search costs are reset and only that network's domain is removed. Previously successful restorations and their occupancy stay in the private transaction. Other networks remain guided. Search failure is an operational trigger, not itself a mathematical infeasibility certificate. A separate capacity diagnosis can justify a stronger statement for a particular fixed scope.

\subsection{Complete checks, acceptance and restoration}
All selected connections must be restored before acceptance. Independent checks validate real edges, unchanged retained branches and terminals, actual global capacities and resource counts. Native FULL and fresh STA evaluate the actual replacement. Acceptance commits the entire group; failure or rejection retains the source state. Intermediate private states may spend resources or slack when a complete mechanism-backed combination is subsequently checked. Requiring every component to improve CPD would exclude some useful groups.

The initial prototype serializes the candidate and performs a separate native LOAD/FULL/STA. A subsequent independent implementation temporarily installs complete candidate copies in the same private worker, runs actual native FULL and a new timing-analysis object, and exports a dedicated candidate route and timing view. It then restores original trees, sink lookups, occupancy, route digest and output hooks, and independently reobserves the source. The worker's original-input timing is never relabeled as candidate timing. This integration has passed small-graph tests and the actual BNN resource transaction: dedicated candidate exports exactly equal a separate same-binary LOAD/FULL/STA, independent physical/scope checks pass, and source recovery is exact. A further same-binary paired export experiment also validates an actual free-path CPD-changing transaction against independent candidate LOAD/FULL/STA. Forced-target-path timing transactions remain outside that validated scope.

Loaded post-legal contexts reconstruct occupancy and native STA but reset the original route run's accumulated PathFinder price history. Matched native/domain controls share this context. An eventual in-router integration must explicitly preserve or reset that history and evaluate the choice.

\section{Overall Development Results}
The implementation derives from VTR revision \texttt{0b912334fdd83bb754b302146ac0d1ec9edc9a0e}. Results below use fixed-input serial clustered routing on the evaluated two-layer, edge-encoded Hybrid-O graphs, ideal clock modeling and native FULL/STA. Flat routing, CHANZ resources and arbitrary multilayer instances are outside this implementation evidence. Each record binds the exact graph, packing, placement, constraints, binary, requests and acceptance rule. A0 is native routing; R1 is the older guided method. Strong BNN A0/R1 roots additionally allow eight native convergence attempts with a1.01 reconvergence threshold. Archived U and Tf predecessors remain separate input/version/convergence records, not interchangeable points in a new algorithm's aggregate score.

\begin{table}[t]
\centering\caption{BNN legal quality points. Fast A0 follows three automatically selected timing commits from pristine strong A0. Resource phases have their own explicitly recorded opportunities. This is a quality comparison, not a matched-compute algorithm ranking.}
\label{tab:bnn}
\small
\begin{tabular}{lrrrr}\hline
State & CPD(ns) & Wire & SB & CB\\\hline
Strong A0 &13.607436&2059567&94&5062\\
Strong R1 &13.421610&2068908&86&5072\\
A0 native K8 &13.607436&2059563&86&5062\\
A0 guided K8 &13.607436&2059842&79&5061\\
Fast A0/native K8 &13.228146&2059603&86&5062\\
Fast A0/guided K8 &13.228146&2059882&79&5061\\
Fast R1/guided K8 &13.228146&2069040&74&5071\\
R1 selective release &13.228146&2069104&72&5071\\\hline
\end{tabular}
\end{table}
The strong-A0 timing trajectory accepts three current-state transactions and reduces CPD by2.787376\%, adding forty wire units with SB/CB unchanged. Its third step improves CPD but regresses a noncritical setup component, making the distinction between strict CPD and endpoint-vector policies observable. The same-source fixed-vertical policy independently selects the identical three states using exact completed-query/transaction reuse; its replay time is not evidence of speedup.

From the fast A0 state, unrestricted K8 resource repair reaches86 SB with wire2059603. Guided K8 reaches79 SB and5061 CB at wire2059882. The additional279 wire units buy seven fewer SB and one fewer CB at exact equal CPD. From the pristine strong A0 root, native K8 instead preserves the slower CPD while reducing wire by four units and SB by eight. The same uniform guided K8 strategy from that pristine root reaches79 SB and5061 CB at2059842 wire, with unchanged13.607436 ns CPD. These slower points use forty fewer wire units than their corresponding timing-first resource endpoints and remain useful under tighter wire requirements. Equal resource shifts do not by themselves prove that the two mechanisms commute; a separately registered resource-first timing phase now tests native CPD repair under the tighter79SB/5061CB cap and the original pristine-root wire envelope. Three automatically accepted transactions now reach13.607436,13.421610,13.413972 and13.228146 ns in order, with wire2059842,2059882,2059878 and2059882;79 SB and5061 CB remain fixed. The last repair expands an observed critical-owner dependency to nine networks and twelve connections. One noncritical connection spends about1.114955 ns of route-delay margin while a critical connection saves about0.185826 ns. Whole native STA improves CPD, but one setup component regresses, so the endpoint-vector policy rejects this step. Relative to the resource-phase input, the complete final vector has eleven improved and seventeen regressed setup components, with no increased hold violation. The fixed source near-critical band has one improvement and no regression. The local current-STA portfolio then yields no further admissible proposal; the research continues with resource and implementation questions.

The resulting four quality coordinates and complete setup/hold endpoint vectors exactly equal the existing timing-first/guided-K8 endpoint, although the serialized route body differs. This is convergence to a known quality point, not a new frontier gain or a general proof that phase order commutes. All ten current winning-connection switch distances equal already-computed exact source-only optimistic references after graph/root/sink identity checks; no queries are repeated for this diagnosis. The reference relaxes occupancy, shared branches and native arithmetic, so it is not a global native-STA bound. These observations favor resource/runtime and other-scope interventions over merely increasing the same cone query budget. Such a candidate would be lost under a universal strict-CPD retention rule.

\begin{table}[t]
\centering\caption{Eltwise own-current K8 development probes from the same pristine native A0 root. The grouping cap and acceptance policy are uniform. Both branches use their own successor states.}
\label{tab:eltwise}
\small
\begin{tabular}{lrrrr}\hline
State & CPD(ns) & Wire & SB & CB\\\hline
Strong A0 root &6.018101&217065&30&968\\
Native K8 &6.018101&217050&26&968\\
Paired guided K8 &6.018101&217214&24&965\\\hline
\end{tabular}
\end{table}
Eltwise guided K8 uses actual group sizes six then three. It saves one additional SB relative to the prior guided K4 endpoint, with identical CPD, wire and CB. No failure release triggers in this circuit, so the gain is attributed to grouping/domain decisions rather than fallback. The native-only branch accepts one move, then rejects three legal but resource-worse alternatives; these rejected candidates remain recorded. On SPMV, paired and native-first R1 repair produce the same legal route at7.045986461 ns,148605 wire,zero SB and967 CB, but remain quality-dominated by strong A0. Initial routability and a circuit-balanced timing improvement across the suite have not been established by these repair probes.

\section{Pareto Tradeoffs and Attribution}
\begin{figure*}[t]
\centering\includegraphics[width=.98\textwidth]{data/joint_quality_projection_03.pdf}
\caption{Full-population and low-resource-detail projections of the complete collected BNN legal population on one physical-input identity, including A0/R1/U/Tf and repairs. Duplicate quality tuples are collapsed. Distinct roots, convergence rules and opportunities remain in the accompanying data. These are not two-dimensional Pareto boundaries: other resource coordinates and complete time are hidden.}
\label{fig:joint}
\end{figure*}
Table~\ref{tab:bnn} and Fig.~\ref{fig:joint} retain external strong baselines. The fast A0 native-K8 point quality-dominates the earlier slower A0 guided-K4 point, which had the same86 SB and5062 CB but wire2059667. Slower native-first low-wire states and older Rq/reassignment states with CB5060 or5059 can still offer distinct tradeoffs. Quality domination is not sufficient for a complete-runtime claim when comparable time is missing.

\subsection{Automatic transfer exposes retained-tree order sensitivity}
A new compact/fused automatic actor from the same pristine eltwise A0 root reaches6.018101129 ns,217050 wire,26 SB and968 CB under unrestricted native repair. Its own-current paired-domain branch instead reaches217191 wire,24 SB and965 CB at the same CPD:141 additional wire units buy two fewer SB and three fewer CB. Both trajectories exhaust their registered current rank-window menu, with one and two accepts respectively. Unlike the old pilot, the new actor has no two-accept local cap. Their complete imported-state workflows are117.15 and118.13 s in this cohort; these are neither cold routing times nor repeated formal timing results. Both use five native children, but fewer child processes do not imply fewer STA evaluations because fused transactions also reobserve restored source. The broader BNN native/paired continuation remains in progress, so historical fixed-prefix points are not declared converged method results.

The23-wire improvement over the older217214-wire guided endpoint has a different cause from reduced graph output. First-step trees and complete endpoint vectors agree, and both second-step request/domain files are identical. Only network6719 changes in the second guided candidate. A same-binary intervention runs the old canonical source, the native-serialized source, and a native source with only6719 restored to canonical branch order. The resulting complete candidate trees and full native endpoint vectors match the old, new and old reference respectively; all pass native FULL/freshSTA and exact source recovery. Retained seed contents agree while their traversal order differs. This establishes search-choice sensitivity to serialization order; it does not prove equal weighted path costs or a particular heap-tie mechanism. Ordered route bytes therefore remain part of search/cache identity. Topology and current timing equivalence alone do not establish equivalent future searches. Confirmation must freeze a traversal policy and give native/domain controls the same policy, rather than crediting this difference to guide scoring or I/O optimization.

\subsection{Capacity feedback changes the domain decision}
An all-eight-domain R1 transaction fails at network31767/pin1 after five successful restorations. A directed path exists when occupancy is ignored. Under the recorded transient occupancy, the retained domain has no capacity path while an unrestricted witness exists. Reconstructing the minimum occupancy immediately after pruning also leaves no domain-capacity path. Because subsequent restorations only add occupancy, changing order or search pressure cannot fix this particular frozen scope/domain. The diagnostic searches a relaxed superset of native re-expansion/closure choices; its negative result is a necessary infeasibility certificate for that scope, not for the entire FPGA.

Releasing only that network's domain yields a fully legal state at72 SB, compared with74 SB for the same-source unrestricted control, at the cost of64 wire units and unchanged CPD/CB. The automatic native release implementation obtains exactly the prior selective-domain route and full endpoint vectors without a manually supplied release network or restarting earlier restorations. Nine setup components improve and none regress, so endpoint-vector acceptance succeeds while strict CPD rejects. All23 components in a fixed source near-critical band stay unchanged. This step demonstrates resource and secondary-timing benefit, not later CPD improvement.

\subsection{Grouping contribution and residual timing potential}
The first native-only fast-A0 K8 group produces exactly the actual route and endpoint vectors of two K4 steps. The first K4 component alone is endpoint-flat and rejected; the complete K8 group improves one setup component with no regressions. This supports complete combination transactions without demanding that every intermediate component pass the final policy. It does not show that a pristine-root endpoint policy reaches that state.

The fast-A0 timing selector eventually has no positive winning-path action. We tested a broader diagnostic: remove all retained-prefix constraints and background occupancy, then query source-root distances for every one of its ten current winning connections. Seven changed queries are newly solved and three exact matches reused. All ten optimistic static-switch distances equal the current values, and every directed witness is independently checked. This limits the evidence for merely increasing the same timing-query budget or rerouting a shared prefix at this state. It is an exact reference for the verified zero-RC switch-delay relaxation, not a full native-arithmetic/STA global optimum or a limit on the overall research scope.

\section{Overhead, Limitations and Next Confirmation}
The automatic selective-release experiment takes287.624 s from its sealed imported source through final context checks. Its native repair core is6.693 s, nested inside a113.480 s worker request. Source materialization takes21.221 s, transaction/domain audits41.032/10.984 s, candidate FULL/STA and materialization65.870 s, final sealing14.856 s, and registration/context/miscellaneous work20.182 s. These stages cannot be added again to the nested native child times. The result motivates state and validation integration; it does not establish lower complete routing time.

The fused candidate implementation passes ten small native jobs, including separate candidate LOAD/FULL/STA equality, unchanged source exports, retained-branch and terminal checks, an infeasible transaction, and failure followed by an exact default request in the same parent. Only a dedicated immutable candidate receipt may materialize candidate timing. An input-state receipt, missing candidate gate or altered route identity is rejected. The BNN actual-candidate physical/materialization checks and separate same-binary native LOAD controls now pass. Two Python comparison/admission errors are preserved: endpoint tuple/list representation and applying a six-digit printed-source CPD check to a full-precision candidate metric. Recovery reuses the sealed native worker and passed transaction/domain audits with zero new repair searches, then executes two new independent validation-only LOAD jobs. No physical or STA tolerance is introduced. Their extra validation-only jobs must not be counted as algorithm work or removed from the reported validation workflow. The new fused kernel COMPLETE interval is28.757 s and includes its candidate observer; connection QUERY intervals sum to7.859 s. The372.048 s recovery segment reuses native creation from a prior attempt and includes new controls, so it is neither a fresh complete algorithm time nor a speedup measurement. A measured I/O target is also apparent: each full timing-graph echo is753.92 MiB and each native frontier357.20 MiB, repeated for source, candidate and restored source. An independent optional graph-text policy now passes ten tiny native jobs and two same-binary BNN transactions, one CPD-changing and one resource-changing. Full and compact modes have identical native searches, actual candidate routes, complete frontier/constraint bytes and source recovery, and dedicated candidates agree with independent LOAD/FULL/STA. Native FULL and new STA still execute in compact mode; it omits only unused full graph text. Timing-cone selection explicitly requires a full export. In the timing case, full versus compact native-child times are81.59 versus62.59 s; complete request times are86.73 versus67.90 s. In the resource case, complete requests are90.21 versus70.66 s. Each compact fused worker avoids about2.21 GiB of graph text while retaining about1.05 GiB of full source/candidate/restored frontier data. Measurement order is reversed between cases. These single paired component observations exclude separately recorded shared startup, use explicit exact-byte physical-audit reuse for the second candidate, and include extra LOAD controls only in validation-workflow totals. They are not formal repeated timing or cold whole-algorithm speedups. The terminal resource-first timing recovery segment takes1793.82 s:894.94 s is query/dependency/portfolio work,834.50 s is nested native transaction coordination, and49.76 s is preparation. The older interrupted segment remains separate; no stages are assembled into a synthetic complete configuration time.

Current evidence has several practical limits. Timing improvements are concentrated in BNN development trajectories; resource transfer does not imply initial routing-success improvement. K1/K4/K8 and scheduling pilots are not one frozen configuration. Complete cold routing-plus-repair time remains unknown, and old U timing, search-core time or sums of historical stages cannot fill it. Confirmation must freeze circuit-independent selection and strategy parameters, start automatically from sealed Tf and strong native roots, and compare A0/R1/native repairs with matched domains, allowed tree scope, target selection and acceptance opportunities. Equal-source local controls provide causal evidence; they do not substitute for that whole-algorithm evaluation.

\subsection{An exact streaming decoder improvement}
The previous endpoint reader retains every per-pin JSON dictionary before extracting the acceptance vector. A new independent reader parses and validates every row but retains only required endpoint/metric data. On two actual357.20 MiB source/candidate traces, old versus new loader times are11.48 versus5.10 s and10.80 versus5.31 s, with reversed measurement order between inputs. Peak process memory falls from1.73 GiB to0.18--0.21 GiB. Complete vector digests, including all52,495 setup components, endpoint/domain identities and full-precision values, are identical. No pin rows are silently skipped; a native COMPLETE pin-count check additionally detects missing rows.

Seven corruption cases are rejected, including malformed pin JSON, missing completion, duplicate or missing endpoints, nonfinite CPD, trailing data and an incomplete pin census. The real already-sealed fused candidate is successfully materialized with identical complete vectors and strict/endpoint decisions, zero new native calls and explicit same-byte physical-audit reuse. These are component memory/latency and integration results under shared server load, not complete routing-plus-repair speedups. Running timing probes retain their registered decoder/version. Complete cold fixed-policy timing remains required.

\section{Conclusion}
Current-STA whole transactions expose useful timing--resource alternatives that strict CPD retention alone misses. Native controls show that selection and grouping already produce gains; retained domains and selective release add different resource tradeoffs under explicit final checks. The next implementation work targets the measured construction/validation cost while preserving actual candidate identity, physical delay, native FULL/STA and source restoration. Broader frozen-policy confirmation and complete runtime measurements remain required. The research goal continues beyond individual probe outcomes.

\begin{thebibliography}{2}
\bibitem{vtr9paper} M. A. Elgammal et al., ``VTR 9: Open-source CAD for fabric and beyond FPGA architecture exploration,'' \emph{ACM Trans. Reconfigurable Technol. Syst.}, 2025, doi:10.1145/3734798.
\bibitem{flatland} C. Hao, A. B. Kahng, B. Pramanik, and I. Youssef, ``Escaping Flatland: A placement flow for enabling 3D FPGAs,'' in \emph{Proc. ACM/IEEE DAC}, 2026, doi:10.1145/3770743.3804424.
\end{thebibliography}
\end{document}
